\chapter{Inspect the stream and accepted file progress}
\MARSChapterHeading{Inspect the stream and accepted file progress}
\label{revision-lab:producing-video-streams}

Read Chapter 27 with this lab. Continue with
\path{revision/labs/files/h263-encoder.asm}; the original \texttt{035-h263}
construction package supports writing the complete encoder in stages. The new
\texttt{h263-read-bits.asm} is an independent unfinished extension. Its feedback
checks field meaning and state, without requiring the worked encoder's labels,
register allocation or legal entropy-code choices.

\section{Trace bits before reading them as bytes}

Write \texttt{A5 A5} as sixteen bits, mark only the first three as valid, then
append the nine-bit value 419. Derive \texttt{BA 35} and committed count twelve.
Mark the four low bits that remain outside the append. Appending four explicit
zeros yields \texttt{BA 30}; simply displaying two bytes did not make those
zeros part of the earlier field.

The \texttt{h263\_write\_bits} context contains pointer, byte capacity and used
bits. Reject an oversized value, a full-width invalid count or an over-capacity
request before buffer access. Test width zero separately: it still validates the
context but reads no buffer and writes no count. Use an output window that begins
inside a guard-filled allocation, so preservation can be checked at both ends.
Repeat every phase from zero through seven, including a 32-bit field crossing
five bytes.

Find the point where the whole append has been admitted. Within that region,
explain the mask and read-modify-write operation for one byte. A correct complete
field cannot be inferred from one correctly written byte: a later space failure
must not leave an earlier prefix published. Larger coefficient and macroblock
stages have separate staging boundaries.

\section{Recover a block's levels from its chosen syntax}

Use INTRA DC code 255, level one at natural index one and minus one at index eight.
The ordinary scan visits those two nonzeros consecutively. Their RUN fields are
zero; LAST is zero then one. Decode both the sixteen-bit compact representation
and the fifty-two-bit escape representation to the same natural-order array.
Do not compare the two outputs byte for byte as a correctness condition.

An INTER nonzero only at the final scan position has RUN 63. Its INTRA AC
counterpart has RUN 62 because the separate DC position is excluded. Move the
last nonzero earlier and verify that LAST moves with it. An all-zero ordinary
scan emits no events; INTRA still carries its DC byte.

Now form presence for six blocks. Six INTRA DC-only blocks have ordinary pattern
zero. For INTER, the four Y presence bits use the complemented CBPY table index.
Keep picture type and macroblock type distinct when selecting MCBPC. Construct
a P-picture INTRA macroblock and explain why it has no motion fields even though
it is inside a predicted picture.

With selected horizontal component minus 32 and predictor 31, recover the wrapped
difference one and the decoder's actual minus 32. Compare coded zero difference
with skipped actual zero motion under a nonzero predictor. Use the full request
contracts when checking these entries; a numerical search result in full pixels
must be converted to half-sample units first.

\section{Run and independently inspect a real clip}

Run the printed helper from a book checkout, using a new directory:
\Needspace{10\baselineskip}
\MARSCodeLabel{Execute and check the complete target encoder}
\begin{Verbatim}
python3 packages/035-h263/construction/h263-clip-construction/run.py \
  revision/labs/files/h263-encoder.asm --output-dir work/run-eight \
  --cli "$CLI" --java "$JAVA"
\end{Verbatim}
The helper needs FFmpeg on the host's command path for independent decoding.
It generates raw QCIF input, retains the prepared source, runs MARS and checks
actual output files. The target program writes the H.263 stream and local
reconstruction. No host encoder fills those outputs for it.

For an IDE observation, use the sibling \texttt{prepare.py} with
\texttt{--source}, \texttt{--output-dir} and the same controls. Open its
\texttt{main.asm} for ordinary MARS Run or Debug. At
\texttt{clip\_caller\_return}, the returned status/count and progress record are
complete, while the caller's named status/count copies are still pending.
A successful default caller prints eight. Retain each experiment's source and
files in its own directory so a later run does not replace the comparison.

The progress record's fields are accepted picture count, accepted encoded byte
count and accepted reconstruction pointer, at native doubleword offsets zero,
eight and sixteen. The reference default produces eight pictures and 2212
encoded bytes. The equivalent escape-oriented source produces 5101 bytes;
that difference alone is not a failure. Inspect the report's picture types,
temporal references, actual motion, skipped macroblocks and per-plane results.

For a separately run checker, point it at the actual files:
\Needspace{11\baselineskip}
\MARSCodeLabel{Check existing files without encoding them again}
\begin{Verbatim}
python3 packages/035-h263/construction/h263-picture-construction/check_stream.py \
  --stream work/run-eight/encoded.263 --input work/run-eight/clip.yuv \
  --reconstruction work/run-eight/reconstructed.yuv \
  --frames 8 --tick-step 6 --report work/stream-report.json
\end{Verbatim}
The report path must be new. The strict parser checks the whole selected profile;
FFmpeg supplies a separate complete decode. Source quality and local/decoder drift
are distinct metrics, and neither replaces the full inverse-accuracy gate.

Vary one controlled choice at a time. Compare scalar and MSA modes with the same
input and coding policy before interpreting work counts. Compare compact and
escape coding without assuming the shorter stream needs fewer instructions.
The 44-picture, six-tick example lasts 8.8088 seconds and ends with temporal
reference two. It checks wraparound without exceeding the documented input limit.

\section{Diagnose failure from the accepted prefix}

Picture publication and file acceptance have different boundaries. A failed
picture retains both caller output areas, although workspace may have changed.
A later file failure can leave a partial output suffix, so inspect its overall
status and the accepted progress record separately.

Force short positive reads and writes and verify that the next operation starts
after the accepted count. A zero or negative write must terminate rather than
spin. The input, stream and trace descriptors can have arbitrary returned values;
only descriptors actually acquired by this call belong in its cleanup set.
The helper's private guard byte is readable scratch, not an additional output byte.

After one accepted picture, reject the next picture's trace write after its stream
write has completed. Progress must still name the preceding reconstruction.
Repeat with a failed stream write after a short positive prefix. Compare physical
file length with accepted encoded bytes; they can differ. Finally append one
trailing input byte. Even when every requested picture was accepted, the exact-EOF
check must reject the whole clip. A count equal to the request does not override
its nonzero status.

Use the existing package domains \texttt{writer}, \texttt{coefficients},
\texttt{syntax}, \texttt{picture} and \texttt{clip} to isolate completed stages.
Their harnesses supply reserved grading storage and documented private stacks.
Do not copy that storage into your program. Keep the original numerical domains
and inverse accuracy checks when changing dependencies of the complete encoder.

\section{Implement the independent bounded reader}

Replace the deliberate \texttt{break} at \texttt{h263\_read\_bits}. Pass an
eight-byte-aligned 32-byte context in \texttt{a0} and the unsigned full-width
requested bit count in \texttt{a1}. Context fields are native doublewords:
input pointer, byte capacity, valid bit count and cursor, at offsets zero,
eight, sixteen and 24.

Capacity is zero through 1048576 bytes. Valid bits are zero through eight times
capacity, and cursor is zero through valid bits. Width is zero through 32.
The complete context and input-capacity extents must not wrap. A positive-capacity
input is mapped and disjoint from context; its byte pointer needs no additional
alignment. Zero capacity never accesses input. These mapped-storage conditions
are caller preconditions, not an instruction to probe arbitrary addresses.

Admit full-width metadata before input reads. Malformed context or width returns
one in \texttt{v0} and zero in \texttt{v1}, retaining every caller byte. A valid
request wider than \texttt{validBits-cursor} returns two and zero, also with no
input access or state change. Subtract only after establishing cursor bounds.
Reading through byte capacity instead of valid bit count could admit padding
as meaningful syntax.

Success returns zero and the right-aligned unsigned field value. A 32-bit result
with its high bit set is still zero-extended in the 64-bit register. Advance only
the cursor, and only after nonempty success. Width zero validates metadata,
returns zero value and performs no input access or context write. Unused bits
after the valid count need not be zero and must not be consumed.

Read only bytes intersecting the selected field. A one-bit read can load its
containing byte, but not an adjacent byte that the field does not touch. Retain
all input bytes and the other three context fields. Preserve saved general
registers, frame pointer, global pointer, stack pointer and the supplied return
continuation. Entry stack alignment is sixteen bytes with 8192 private bytes;
do no console or file I/O.

One approach extracts a single bit each iteration. Another consumes the available
fragment of each byte and shifts the accumulator by that fragment's width.
Both can preflight the complete requested width and commit the final cursor once.
Do not carry the previous result into a new call's accumulator. A success followed
by a rejected read must return zero value for the rejection.

\section{Use the reader to explain an actual header}

The small demo supplies bytes \texttt{BA 35}, valid count twelve and cursor three.
At \texttt{reader\_demo\_return}, a completed function returns zero and 419,
with cursor twelve. Continuing copies the return values and exits without console
output. Its next one-bit read must fail; a zero-width read must succeed without
changing cursor. Also start at cursor zero: widths three then nine return five
and 419, respectively.

Run its separate feedback from the book root:
\Needspace{8\baselineskip}
\MARSCodeLabel{Check field values, access bounds and cursor ownership}
\begin{Verbatim}
java -jar "$CLI" grade \
  revision/labs/checks/h263-read-bits.kts \
  revision/labs/files/h263-read-bits.asm
\end{Verbatim}
The unfinished body is expected to fail until implemented. Checks cover every
width/byte phase, both byte orders, unsigned full fields, maximum capacity,
invalid full-width metadata, exact touched bytes, empty reads and reused state.
Equivalent bit and byte-fragment traversals are accepted.

For an extension, copy the first seven encoded bytes of a real picture and give
the reader fifty valid bits. Read widths 22, 8, 13, 5, 1 and 1. The first field
is numerical start code 32; the next is its temporal reference. The remaining
fields describe the restricted picture type, quantizer and disabled flags.
Compare all eight pictures, then the wrapped references in the longer clip.
This inspects a header; it does not parse the intervening variable-length block
payloads or reconstruct pixels. Explain exactly what your new inspector proves
before using it as a decoder check.
